\documentclass[10pt,a4paper]{amsart}
\usepackage[margin=19mm]{geometry}
\usepackage{amsmath,amssymb,mathtools}
\usepackage[scr=rsfs,cal=euler]{mathalfa}
\usepackage{xfrac}
\usepackage[unicode,hidelinks,bookmarks=false]{hyperref}
\newcommand{\ie}{i.e.\ }
\newcommand{\cf}{cf.\ }
\newcommand{\resp}{respectively\ }
\newcommand{\Subsection}{\subsection}
\providecommand{\nobreakdash}{\nobreak\hbox{-}}
\setlength{\emergencystretch}{2em}
\multlinegap0pt
\usepackage{mathtools}
\usepackage[shortlabels]{enumitem}
\usepackage{amssymb}
\usepackage{float}
\usepackage{algorithm}\usepackage{algpseudocode}
\newtheorem{lemma}{Lemma}
\newtheorem{claim}{Claim}
\newtheorem{prop}{Proposition}
\usepackage{cleveref}
\crefname{lemma}{Lemma}{Lemmas}
\crefname{claim}{Claim}{Claims}
\crefname{prop}{Proposition}{Propositions}
\newcommand{\cH}{\mathcal{H}}
\newcommand{\eps}{\varepsilon}
\renewcommand{\setminus}{-}
\title{Nearly Hamilton cycles in sublinear expanders, and applications}
\author{Shoham Letzter, Abhishek Methuku, Benny Sudakov}
\date{Source: arXiv:2503.07147v2 (CC BY 4.0)}
\begin{document}
\maketitle

\noindent\emph{Source and licence.} Nearly Hamilton cycles in sublinear expanders, and applications. Version arXiv:2503.07147v2 (CC BY 4.0), distributed by arXiv under the Creative Commons Attribution 4.0 International licence, http://creativecommons.org/licenses/by/4.0/, as recorded in the arXiv licence metadata for this exact version and shown on the abstract page https://arxiv.org/abs/2503.07147v2. Full licence notice: \texttt{licenses/arxiv-2503.07147-CC-BY-4.0.txt}.\par\smallskip

\noindent\textit{Editorial scope.} Counted unit: the article's lemma lem:expandercoverweak (source0.tex lines 478--481). The article's complete original proof of it is delivered with it (source0.tex lines 555--666, one \texttt{proof} environment of 112 lines, which the article places away from the statement). No mathematics is written, repaired or re-derived: the statement and the proof are the article's own lines, in the article's own notation, and no retained line is substituted or re-flowed.  Retained uncounted as the source-local context the proof needs: lines 526--529.  Nothing was omitted from any retained span, so the package carries no omission record.  No retained line contains a citation, so the wrapper carries no bibliography and no bibliography entry.  No retained line contains a figure environment, an image-inclusion command or drawing code, so no image is delivered and the package declares no asset dependency.  Editorial formatting only, in addition: an AMS \texttt{amsart} preamble that restates the article's own declarations and macros used by the retained lines, namely the environments \texttt{lemma}, \texttt{claim}, \texttt{prop} on separate counters, together with the standard packages those lines need.  The retained environments print in this document as Lemma 1, Claim 1, Proposition 1 in order of appearance, whereas the article prints the same items with the numbering of its own full document.  The complete native source of the article as distributed in the arXiv e-print for this exact version is bundled unchanged as \texttt{sources/174484/source0.tex}.
\begin{lemma}
	\label{lem:packingoneround}
	Let $\eps, \lambda, n$ satisfy $\lambda \log n \le \min \{\frac{1}{10}, \eps \}$. Let $G$ be an $n$-vertex graph with $\Delta(G) \le d$ and $d(G) \ge d(1- \eps)$. Then, there exists a collection $\cH$ of vertex-disjoint $\lambda$-expanders in $G$ such that $d(H) \ge d(1-10 \eps)$ and $\delta(H) \ge d(H)/2$ for every $H \in \cH$ and $\sum_{H \in \cH} |V(H)| \ge n/4$.
\end{lemma}

\begin{lemma}
	\label{lem:expandercoverweak}
	Let $\alpha > 0$ be a fixed real number and let $\eps, \lambda, n$ satisfy $\lambda \log n \le \min \{\frac{1}{10}, \eps \}$ and let $n$ be sufficiently large. Let $G$ be an $n$-vertex graph with $\Delta(G) \le d$ and $d(G) \ge d(1-\eps)$. Then, there exists a collection $\cH$ of vertex-disjoint $\lambda$-expanders in $G$ such that $d(H) \ge d(1 - \eps (\log n)^{28 \alpha})$ and $\delta(H) \ge d(H)/2$ for every $H \in \cH$ and $\sum_{H \in \cH} |V(H)| \ge (1 - \frac{1}{(\log n)^{\alpha}}) n$. 
\end{lemma}

\begin{proof}[ of \Cref{lem:expandercoverweak}]
	Let $C = 100$. For every $i \ge 1$, let $\eps_i = C^{i} \eps$, and let $\cH_i$ be a collection of vertex-disjoint $\lambda$-expanders with average degree at least $d(1- \eps_i)$ in $G \setminus (\bigcup_{H \in \cH_1 \cup \ldots \cup \cH_{i-1}} V(H))$ that maximises the number of vertices covered by expanders in $\cH_i$ among all such collections. For all $i \ge 1$, let $V_i \coloneqq \bigcup_{H \in \cH_i} V(H)$. 

	\begin{claim}
		\label{claim:coverbyexpanders}
		For all $i \ge 1$, $\sum_{k = 1}^i |V_k| \ge (1 - (\frac{3}{4})^{i})n$.
	\end{claim}
	\begin{proof}[ of claim] 
		We prove the claim by induction. For $i = 1$, the claim immediately follows by applying Lemma~\ref{lem:packingoneround}
		to $G$ to obtain a collection $\cH_1$ of vertex-disjoint $\lambda$-expanders covering at least $n/4$ vertices and satisfying $d(H) \ge d(1 - 10 \eps) \ge d(1 - C \eps) = d (1 - \eps_1)$ and $\delta(H) \ge d(H)/2$ for every $H \in \cH_1$. This shows that $|V_1| \ge n/4$, as desired.

		Now suppose that for all $1 \le j \le i$, we have $\sum_{k = 1}^j |V_k| \ge (1 - (\frac{3}{4})^j)n$. We will show that $\sum_{k = 1}^{i+1} |V_k| \ge (1 - (\frac{3}{4})^{i+1})n$. 

		Let $U \coloneqq V_1 \cup \ldots \cup V_i$ and let $\overline{U} \coloneqq V(G) \setminus U$. Since for each $1 \le k \le i$, the average degree of the $\lambda$-expanders in $\cH_k$ (whose union spans the vertex set $V_k$) is at least $d(1- \eps_k)$, we have $e(U) \ge \sum_{k=1}^i \frac{d}{2} (1- \eps_k) |V_k| = \frac{d}{2} |U| - \frac{d}{2} \sum_{k=1}^i \eps_k|V_k|$. Hence, 
		\begin{equation}
			\label{eq:claim2:1}
			e(U) + e(U, \overline{U}) \le d |U| - e(U) \le \frac{d}{2} |U| + \frac{d}{2} \sum_{k=1}^i \eps_k|V_k|.
		\end{equation}

Note that since $\eps_1 \le \ldots \le \eps_i$ and $\sum_{k = 1}^j |V_k| \ge (1 - \left(\frac{3}{4}\right)^j)n$ for $1 \le j \le i$, we can bound the sum $\sum_{k=1}^i \eps_k|V_k|$ using the proposition below. 

\begin{prop}\label{prop:weighted-sum-bound}
Let $\eps_1\le \cdots \le \eps_i$ and let $a_1,\ldots,a_i\ge 0$ satisfy $\sum_{k=1}^i a_k \le n$ and
$\sum_{k=1}^j a_k \ge \big(1-(\tfrac34)^j\big)n$ for all $1\le j\le i$.
Then,
\[
\sum_{k=1}^i \eps_k a_k
\le
\sum_{k=1}^{i-1}
\eps_k(\tfrac34)^{k-1}\tfrac n4
+\eps_i(\tfrac34)^{i-1}n .
\]
\end{prop}

\begin{proof}
Note that since $\eps_1\le\cdots\le\eps_i$, the sum 
$\sum_{k=1}^i \eps_k a_k$ 
is maximised when, for each $1\le k\le i$, the value of $a_k$ is taken to be as large 
as permitted by the inequalities 
$\sum_{j=1}^k a_j \ge (1-(\tfrac34)^k)n$ and $\sum_{k=1}^ia_k \le n$, 
assuming that the values of $a_{k+1},\ldots,a_i$ have already been fixed at their 
largest admissible values.

From $\sum_{k=1}^{i-1} a_k \ge (1-(\tfrac34)^{i-1})n$ and $\sum_{k=1}^i a_k \le n$, we obtain
$a_i \le (\tfrac34)^{i-1}n$. By choosing $a_i$ as large as possible, we obtain $a_i = (\tfrac34)^{i-1}n$, which implies $\sum_{k=1}^{i-1}a_k = (1 - (\tfrac34)^{i-1})n$.
Similarly, using
$\sum_{k=1}^{i-2} a_k \ge (1-(\tfrac34)^{i-2})n$ we get
\[
a_{i-1}
\le
\big((\tfrac34)^{i-2}-(\tfrac34)^{i-1}\big)n
=
(\tfrac34)^{i-2}\tfrac n4 .
\]

Again, by choosing $a_{i-1}$ as large as possible, we obtain $a_{i-1} =(\tfrac34)^{i-2}\tfrac n4$. Proceeding inductively, for every $1\le k\le i-1$, we obtain
$a_k = (\tfrac34)^{k-1}\tfrac n4$.
Substituting these values of $a_k$ for $1 \le k \le i-1$, and $a_i$ yields
\[
\sum_{k=1}^i \eps_k a_k
\le
\sum_{k=1}^{i-1}
\eps_k(\tfrac34)^{k-1}\tfrac n4
+\eps_i(\tfrac34)^{i-1}n ,
\]
as required.
\end{proof}
By Proposition~\ref{prop:weighted-sum-bound} (applied with $a_k=|V_k|$), we have
%\[
%\sum_{k=1}^i \eps_k|V_k|
%\le \sum_{k=1}^{i-1} %\eps_k(\tfrac34)^{k-1}\tfrac n4 + \eps_i(\tfrac34)^{i-1}n.
%\]
%Note that since $\eps_1 \le \ldots \le \eps_i$ and $\sum_{k = 1}^j |V_k| \ge (1 - \left(\frac{3}{4}\right)^j)n$ for $1 \le j \le i$, the sum $\sum_{k=1}^i \eps_k|V_k|$ is maximised if, for each $1 \le k \le i$, the size of $V_k$ is chosen to be as large as possible, assuming the sizes of $V_{k+1}, \ldots, V_i$ were already chosen to be as large as possible. Since $\sum_{k = 1}^{i-1} |V_k| \ge (1 - (\frac{3}{4})^{i-1})n$, the largest possible size of $V_i$ is $(\frac{3}{4})^{i-1} n$. Given this, the largest possible size of $V_{i-1}$ is $(1 - (\frac{3}{4})^{i-1})n - (1 - (\frac{3}{4})^{i-2})n = ((\frac{3}{4})^{i-2} - (\frac{3}{4})^{i-1})n = (\frac{3}{4})^{i-2} \frac{n}{4}$, using $\sum_{j = 1}^{i-2} |V_j| \ge (1 - (\frac{3}{4})^{i-2})n$. Similarly, for each $1 \le k \le i-1$, the largest possible size of $V_k$ is $(\frac{3}{4})^{k-1} \frac{n}{4}$. Hence, 
		\begin{align}
			\label{eq:claim2:2}
			\begin{split}
				\sum_{k=1}^i \eps_k|V_k| 
				& \le \sum_{k=1}^{i-1} \eps_k \left(\frac{3}{4}\right)^{k-1} \frac{n}{4} + \eps_i \left(\frac{3}{4}\right)^{i-1} n  \\
				& \le \sum_{k=1}^{i} \frac{4}{3} \eps \left(\frac{3 C}{4}\right)^{k} n 
				\le \frac{\frac{4}{3} \eps}{\frac{3 C}{4}-1} \left(\frac{3 C}{4}\right)^{i+1}n.
			\end{split}
		\end{align}
		Combining \eqref{eq:claim2:1} and \eqref{eq:claim2:2}, we have $$e(U) + e(U, \overline{U}) \le \frac{d}{2} |U| + \frac{\frac{2}{3} \eps d}{\frac{3 C}{4}-1} \left(\frac{3 C}{4}\right)^{i+1}n.$$
		Therefore, 
		\begin{align} \label{claim2:eq3}
			\begin{split}
				e(\overline{U}) \ge e(G) - (e(U) + e(U, \overline{U})) 
				&\ge \frac{1}{2} nd(1-\eps) - \frac{d}{2} |U| - \frac{\frac{2}{3} \eps d}{\frac{3 C}{4}-1} \left(\frac{3 C}{4}\right)^{i+1}n \\  
				&\ge \frac{d}{2} |\overline{U}| - \frac{1}{2} n d \eps - \frac{8\eps d}{9C - 12} \left(\frac{3 C}{4}\right)^{i+1}n. 
			\end{split}
		\end{align}

		Since $\sum_{k = 1}^i |V_k| \ge (1 - (\frac{3}{4})^i)n$, we have $|\overline{U}| \le (\frac{3}{4})^i n$. Thus, by \eqref{claim2:eq3}, 
		\begin{align*}
			d(\overline{U}) 
			= \frac{2 e(\overline{U})}{|\overline{U}|} 
			& \ge d\left(1 - \left(\frac{4}{3}\right)^i \eps \left(1 + \frac{16}{9C - 12} \left(\frac{3 C}{4}\right)^{i+1}  \right) \right) \\
			& = d \left(1 - \eps \left(\left(\frac{4}{3}\right)^i + \frac{4 C^{i+1}}{3C - 4}\right) \right).
		\end{align*}

		Note that since $C = 100$, we have $\left(\frac{4}{3}\right)^i + \frac{4 C^{i+1}}{3C - 4} \le \frac{C^{i+1}}{10}$, so
		\begin{equation*}
			d(\overline{U}) \ge d \left(1 - \frac{\eps C^{i+1}}{10}\right)  = d \left(1 - \frac{\eps_{i+1}}{10}\right).
		\end{equation*}
		Now, note that $\lambda \log |\overline{U}| \le \lambda \log n \le \min \{ \frac{1}{10},  \eps \} \le \min \{ \frac{1}{10},  \frac{C^{i+1} \eps}{10} \} = \min \{ \frac{1}{10},  \frac{\eps_{i+1}}{10} \}$, where in the last inequality we used $C = 100$. Hence, by Lemma~\ref{lem:packingoneround}, there exists a collection $\cH_{i+1}$ of vertex-disjoint $\lambda$-expanders in $G[\overline{U}]$ with $d(H) \ge d(1 - \eps_{i+1})$, and $\delta(H) \ge d(H)/2$ for every $H \in \cH_{i+1}$ such that, if $V_{i+1} = \bigcup_{H \in \cH_{i+1}} V(H)$, then $|V_{i+1}| = |\bigcup_{H \in \cH_{i+1}} V(H)|\ge |\overline{U}|/4$. Therefore, $n - \sum_{k = 1}^{i+1} |V_k| \le \frac{3}{4}|\overline{U}|$. Since $|\overline{U}| \le (\frac{3}{4})^i n$, we have 
		$ n - \sum_{k = 1}^{i+1} |V_k| \le \frac{3}{4}|\overline{U}| \le (\frac{3}{4})^{i+1}n.$ This shows that $\sum_{k = 1}^{i+1} |V_k| \ge (1-(\frac{3}{4})^{i+1})n$, proving the claim.
	\end{proof}

Let $t = \min \{ i \mid  (\frac{3}{4})^{i+1} \le \frac{1}{(\log n)^{\alpha}}\}$. Then, $(\frac{3}{4})^t \ge \frac{1}{(\log n)^{\alpha}}$, so $(\frac{4}{3})^t \le (\log n)^{\alpha}$ which implies that $t \le \frac{\alpha }{\log (4/3)} \log\log n \le 3 \alpha \log\log n$. Thus, $C^{t+1} \le C^{3 \alpha \log\log n + 1} \le C^{4 \alpha \log\log n} \le 2^{28 \alpha \log\log n}  = (\log n)^{28 \alpha}.$ 

Now consider the collection of $\lambda$-expanders $\cH \coloneqq \cH_1 \cup \ldots \cup \cH_{t+1}$. Every $H \in \cH$ is a $\lambda$-expander with $d(H) \ge d(1 - \eps_{t+1}) = d(1 - \eps C^{t+1}) \ge d(1 - \eps (\log n)^{28 \alpha})$, $\delta(H) \ge d(H)/2$, and by Claim~\ref{claim:coverbyexpanders}, these expanders cover at least $\sum_{k = 1}^{t+1} |V_k| \ge (1 - (\frac{3}{4})^{t+1})n \ge (1 - \frac{1}{(\log n)^{\alpha}})n$ vertices of $G$, proving the lemma.
\end{proof}

\end{document}
